\documentclass[conference]{IEEEtran}
\usepackage{cite}
\usepackage{amsmath,amssymb,amsfonts}
\usepackage{algorithmic}
\usepackage{graphicx}
\usepackage{textcomp}
\usepackage{xcolor}
\usepackage{booktabs}

\begin{document}

\title{Deterministic Semantic Guardrails for Autonomous LLM-Driven Site Reliability Engineering in Cloud-Native Systems}

\author{\IEEEauthorblockN{Tanya}
\IEEEauthorblockA{\textit{Cloud and Systems Engineering} \\
Repository: https://github.com/Tanyarajput677/ai-sre-agent}
}

\maketitle

\begin{abstract}
Autonomous incident remediation using Large Language Models (LLMs) in distributed Kubernetes environments introduces non-trivial operational risks, including hallucinated configurations, privilege escalation, and destructive cluster mutations. We present an autonomous AI-SRE framework that unifies in-cluster telemetry extraction, automated root-cause analysis (RCA), and a deterministic semantic guardrail layer. Our system enforces declarative GitOps-compliant state transitions while blocking arbitrary execution paths. Benchmarking across common Kubernetes failure signatures (CrashLoopBackOff, ImagePullBackOff, and probe misconfigurations) demonstrates sub-second Mean Time to Remediation (MTTR < 0.05s) with zero policy violations.
\end{abstract}

\begin{IEEEkeywords}
Site Reliability Engineering, Kubernetes, Large Language Models, Guardrails, Autonomous Remediation, Cloud-Native, GitOps.
\end{IEEEkeywords}

\section{Introduction}
Modern microservice architectures require sub-minute incident response times to meet Service Level Objectives (SLOs). Traditional human-in-the-loop SRE workflows incur substantial triage latency (MTTR $> 2-5$ minutes). While generative AI agents show capability in root-cause isolation, executing LLM-generated operations directly against a cluster introduces security risks:
\begin{enumerate}
    \item Execution of destructive commands (\texttt{rm -rf}, shell injection).
    \item Unauthorized privilege escalations (\texttt{privileged: true}, \texttt{hostPID}).
    \item Non-declarative imperative mutations violating GitOps principles.
\end{enumerate}

\section{System Architecture}
The framework operates as a closed-loop control system comprising four layers:
\begin{itemize}
    \item \textbf{Telemetry Ingestion Engine:} Directly interfaces with the Kubernetes CoreV1 API to extract container states, exit codes, and localized event streams without sidecar overhead.
    \item \textbf{Diagnostic Reasoning Loop:} Generates structured Root Cause Analysis (RCA) metadata and declarative patch proposals.
    \item \textbf{Deterministic Semantic Guardrails:} An AST-based policy engine enforcing security constraints and YAML schema validity prior to reconciliation.
    \item \textbf{Declarative Reconciliation:} Applies verified patches idempotently via the Kubernetes control plane.
\end{itemize}

\section{Experimental Evaluation}
We evaluated the system on a multi-fault KinD testbed across three failure modes:
\begin{table}[htbp]
\caption{Benchmark Latency and Remediation Performance}
\begin{center}
\begin{tabular}{@{}lccc@{}}
\toprule
\textbf{Failure Scenario} & \textbf{Telemetry Latency} & \textbf{Remediation MTTR} & \textbf{Guardrail Status} \\ \midrule
CrashLoopBackOff          & 0.018s                     & 0.035s                    & PASSED (100\%)            \\
Probe Failure             & 0.019s                     & 0.042s                    & PASSED (100\%)            \\
ImagePullBackOff          & 0.017s                     & 0.038s                    & PASSED (100\%)            \\ \bottomrule
\end{tabular}
\end{center}
\end{table}

\section{Conclusion}
By decoupling diagnostic reasoning from deterministic policy verification, autonomous AI-SRE frameworks can safely achieve sub-second MTTR in production-grade Kubernetes topologies.

\end{document}
